% ===================================================================
%  TABLEAU N° 1 — L'école, le lycée, la salle de classe.
%
%  Ceci n'est PAS une transcription : c'est une traduction, faite en
%  2026, en francais standard du Quebec. D'ou trois differences avec
%  texto/io et texto/fr, et trois seulement :
%
%   * pas de \nl ni de \cc — la traduction ne suit aucune ligne
%     imprimee, et les lignes de ce fichier ne sont que du confort ;
%   * pas de \begin{VUpage} — elle n'a ni feuillet ni folio, et la
%     page de lecture ne lui en affiche pas ;
%   * le gras se pose sur le TERME seul, ponctuation dehors.
%
%  Tout le reste est identique, et doit l'etre : les cles « %%K » sont
%  celles de texto/io, macro pour macro pour l'apparat, et les renvois
%  \textsuperscript{(n)} visent les MEMES objets de la planche — ce
%  sont eux qui ouvrent les gros plans.
%
%  CETTE COLONNE SE TRADUIT SUR L'IDO, PAS SUR LE FAC-SIMILE FRANCAIS.
%  C'etait le point a trancher, et il a ete tranche : une colonne
%  fr-CA calquee sur texto/fr aurait ete une VARIANTE — le dispositif
%  de texto/varianti.json, six paires deja en place (en-GB/en-US,
%  es-ES/es-MX, de-DE/de-AT, pt-PT/pt-BR, nl-NL/nl-BE, zh-Hans/
%  zh-Hant) — et une variante se derive par remplacements sur sa base.
%  Or texto/fr est un fac-simile de 1926 : figee, elle n'est pas une
%  base sur laquelle on remplace, elle est une source qui ne bouge
%  pas. Deriver le quebecois de 2026 d'un francais de 1926 aurait
%  demande de reecrire la phrase entiere a chaque bloc, ce qui n'est
%  plus un calque mais une traduction qui se cache. Elle se fait donc
%  au grand jour : trente-et-unieme colonne, seize fichiers, traduite
%  de l'ido comme les trente autres.
%
%  D'OU CE QUE LA COLONNE N'EST PAS : elle n'est pas texto/fr rendu
%  quebecois. Les deux colonnes traduisent le meme ido et se
%  ressemblent donc beaucoup — l'apparat du volume y est mot pour mot
%  le meme, et il le faut, puisque c'est la meme langue qui traduit le
%  meme titre. Ce qui les separe se compte : sur ce tableau-ci, neuf
%  choix lexicaux et une convention typographique.
%
%  LES NEUF CHOIX, ET POURQUOI. Aucun n'est un regionalisme jete pour
%  la couleur : chacun est le mot que l'ecole quebecoise emploie
%  aujourd'hui, la ou le livret de 1926 en emploie un autre.
%
%    (36) gomme a effacer ....... efface
%    (57) sac de classe ......... sac d'ecole
%    (37) porte-crayon .......... etui a crayons
%    (65) pupitre ............... lutrin  [le pupitre a musique]
%    (74) repetiteur ............ surveillant
%    (72) censeur ............... directeur des etudes
%    (58) portemanteaux ......... crochets
%         interne ............... pensionnaire
%         laborieux ............. travaillant
%
%  Et une dixieme ligne qui n'est pas un choix lexical : l'ido ecrit
%  « docas », enseigner, non « aprender » — le professeur ENSEIGNE aux
%  eleves a parler, il ne leur APPREND pas a parler.
%
%  « CAR » N'EST PAS EN CAUSE, et une premiere redaction de cette
%  liste l'y avait mis a tort. Le mot n'est ni vieilli ni etranger au
%  Quebec : il s'ecrit en 2026 comme il s'ecrivait en 1926, des deux
%  cotes de l'Atlantique. Ce qui l'a fait remplacer une fois, au bloc
%  t01-11-1, n'est pas son registre mais son voisinage : la phrase y
%  portait deja deux points-virgules, et « parce qu' » y coupe mieux.
%  Ailleurs dans la colonne, « car » reste.
%
%  ET DEUX PIEGES EVITES, qui sont l'inverse : des mots ou le Quebec
%  dit AUTRE CHOSE que ce que le lecteur de France attendrait, et ou
%  la traduction doit donc s'ecarter du mot quebecois.
%    * « cartable » au Quebec est un classeur a anneaux, non un sac :
%      le sac (57) est un « sac d'ecole ».
%    * « pupitre » au Quebec est le bureau de l'eleve. Mais les tables
%      (18) de cette classe portent TROIS eleves chacune — la gravure
%      le montre et le texte le dit — ce ne sont donc pas des
%      pupitres : ce sont des tables. Le mot « pupitre » reste libre,
%      et sert la ou il ne trompe pas : (65), le pupitre a musique,
%      qui se dit « lutrin » ici.
%
%  LA CONVENTION TYPOGRAPHIQUE : pas d'espace devant « ; », « ! » et
%  « ? », espace devant « : ». C'est l'usage courant du Quebec, ou
%  l'espace fine insecable recommandee devant les trois premiers est,
%  faute d'etre disponible, ordinairement omise plutot que remplacee
%  par une espace pleine. Le fac-simile de 1926, lui, la pose pleine
%  (« langues vivantes ? ») : c'est la, en une marque, la difference
%  qui se voit le plus vite entre les deux colonnes francaises.
%
%  ON MODERNISE LA LANGUE, JAMAIS LES CHOSES. Le lycee reste un lycee,
%  l'encrier un encrier, l'ardoise une ardoise, le poele un poele : ce
%  sont les objets de la planche, et la planche ne change pas. Ce qui
%  se modernise, c'est le mot par lequel un lecteur d'aujourd'hui les
%  nommerait.
%
%  LES PRENOMS SONT FRANCAIS, comme dans texto/fr : Ioannes est Jean,
%  Iakobus est Jacques, Karolus est Charles. Les noms latins de la
%  note (*) ne bougent pas : c'est d'eux que la note parle.
%
%  « postigo » (t01, renvoi 78) EST UN VASISTAS. Verifie sur la
%  gravure : la verriere de la classe a deux rangees de carreaux, et
%  dans la rangee HAUTE deux petits vantaux sont graves battants,
%  gonds sur un montant vertical, ouverts vers l'interieur. C'est le
%  seul ouvrant de la verriere, et il ne sert qu'a aerer. Le Quebec
%  n'a pas de mot a lui pour cet objet — « chassis » y nomme la
%  fenetre entiere, ce qui serait un contresens — et garde donc
%  « vasistas », que les dictionnaires quebecois enregistrent.
% ===================================================================

%%K t01-apar-0 apar
\VUsaut{2.0mm}
\VUcentre{12.6pt}{120}{LIVRET EXPLICATIF}
\VUsaut{3.2mm}
\VUcentre{8.4pt}{160}{DES}
\VUsaut{2.6mm}
\VUtitre{15.6pt}{0}{Tableaux Auxiliaires Delmas}
\VUsaut{3.4mm}
\VUornamento{9pt}
\VUsaut{5.0mm}
\VUcentre{13.2pt}{90}{PREMIÈRE SÉRIE}
\VUsaut{3.0mm}
\VUfilet{16mm}
\VUsaut{5.6mm}

%%K t01-apar-1 apar
\VUtitre{15.0pt}{60}{TABLEAU N\textsuperscript{o} 1}
\VUsaut{1.4mm}
\VUtitre{11.4pt}{0}{L'école, le lycée, la salle de classe.}
\VUsaut{2.2mm}
\VUfilet{20mm}
\VUsaut{6.4mm}

%%K t01-tit-1 sub
\VUpk{10.2pt}{40}{Description générale.}
\VUsaut{3.0mm}

%%K t01-01-1 p
1. --- Ce tableau représente une \VUgras{salle de classe} dans un
\VUgras{lycée}. Dans cette salle, il y a huit \VUgras{tables}
\textsuperscript{(18)} et huit \VUgras{bancs} \textsuperscript{(32)}. Trois
élèves sont assis sur chaque banc. Le \VUgras{professeur}
\textsuperscript{(7)} est assis sur une \VUgras{chaise}
\textsuperscript{(8)} dans sa \VUgras{chaire} \textsuperscript{(6)}.

%%K t01-02-1 p
2. --- À gauche de la chaire se trouve une \VUgras{armoire vitrée}
\textsuperscript{(15)}. À droite, il y a le \VUgras{tableau noir}
\textsuperscript{(1)} avec la \VUgras{brosse} \textsuperscript{(2)} et la
\VUgras{craie} \textsuperscript{(3)}.

%%K t01-02-2 p
Au-dessus de la chaire sont accrochés des \VUgras{tableaux muraux}
\textsuperscript{(9, 11, 12)} et une \VUgras{carte} \textsuperscript{(10)}.

%%K t01-03-1 p
3. --- Les élèves ne sont pas tous à leur \VUgras{place}
\textsuperscript{(17)}. Jean \textsuperscript{(4)} est debout devant le
tableau; il tient la brosse et efface les \VUgras{figures de géométrie}.

%%K t01-03-2 p
Pierre est près de la chaire; il ramasse les \VUgras{devoirs écrits}
\textsuperscript{(5)} et les apporte au professeur.

%%K t01-04-1 p
4. --- \VUgras{Ce} professeur est le professeur de \VUgras{langues
vivantes}. Il enseigne aux élèves à parler, à lire et à écrire des langues
étrangères \VUgras{(allemand, anglais, espagnol, italien, russe} ou
\VUgras{français)}.

%%K t01-05-1 p
5. --- La salle de classe est très grande et très claire. Au fond, il y a une
large \VUgras{fenêtre} munie de deux \VUgras{vasistas}
\textsuperscript{(78)} et une \VUgras{porte vitrée} pour l'aérer et
l'éclairer.

%%K t01-05-2 p
Cette salle de classe n'est pas froide. Au milieu, pour la chauffer, il y a un
très bon \VUgras{poêle} \textsuperscript{(46)} avec son \VUgras{tuyau}
\textsuperscript{(45)}.

%%K t01-05-3 p
Des \VUgras{crochets} \textsuperscript{(58)} sont fixés au mur; les
casquettes et les manteaux des élèves y sont suspendus.

%%K t01-tit-2 sub
\VUsaut{5.2mm}
\VUpk{10.2pt}{40}{La cour de récréation.}
\VUsaut{3.6mm}

%%K t01-06-1 p
6. --- L'\VUgras{horloge} \textsuperscript{(63)} marque neuf heures cinq.

%%K t01-06-2 p
La \VUgras{récréation} \textsuperscript{(76)} vient de finir et la leçon
recommence. La récréation se prend dans la \VUgras{cour}
\textsuperscript{(75)}. On voit quelques élèves qui jouent. Un
\VUgras{surveillant} \textsuperscript{(74)} veille sur eux.

%%K t01-07-1 p
7. --- Dans la cour, on voit encore d'autres personnes : l'\VUgras{inspecteur}
\textsuperscript{(73)}, le chef d'établissement \VUgras{(le directeur)}
\textsuperscript{(71)} et le \VUgras{directeur des études}
\textsuperscript{(72)}.

%%K t01-07-2 p
L'inspecteur inspecte les écoles, le directeur dirige le lycée, le directeur
des études surveille les études.

%%K t01-07-3 p
Le quatrième personnage est le \VUgras{concierge} \textsuperscript{(77)}. Il
porte sous le bras un registre pour y inscrire les absences.

%%K t01-08-1 p
8. --- Au fond de la cour se trouvent les \VUgras{appareils de gymnastique}
\textsuperscript{(70)} et l'atelier des \VUgras{travaux manuels}
\textsuperscript{(69)}.

%%K t01-08-2 p
À droite du tableau, on commence à voir les \VUgras{salles} de
\VUgras{musique} et de \VUgras{dessin}.

%%K t01-noto-1 noto
\VUnotes{143.2mm}{%
(*) Pour les \textit{prénoms}, on conseille de les transcrire aussi sous leur
forme latine. C'est le seul moyen d'échapper à d'innombrables variantes.
D'ailleurs, comment choisir entre \textit{Giovanni, Jean,} \textit{Johann, Jan,
Hans, John, Ivan,} etc.? Allons-nous créer un dictionnaire spécial des prénoms?
Puisons dans le latin leur nominatif et disons : \VUgras{Ioannes, Ioanna;
Iulius, Iulia; Iakobus; Andreas; Lukas.} (Grammaire complète).%
}

%%K t01-tit-3 sub
\VUpk{10.2pt}{40}{Description détaillée.}
\VUsaut{2.4mm}

%%K t01-tit-4 sub
\VUpk{10.2pt}{40}{Les bons élèves.}
\VUsaut{3.0mm}

%%K t01-09-1 p
9. --- Les élèves du premier banc à droite de la chaire sont \VUgras{Henri}
(*), \VUgras{Étienne} et \VUgras{Charles}.

%%K t01-09-2 p
Henri est très attentif et travaillant; il écoute le professeur. Sur sa table,
il a un \VUgras{cahier} \textsuperscript{(28)}, un \VUgras{livre}
\textsuperscript{(29)}, un \VUgras{dictionnaire} \textsuperscript{(30)} et un
\VUgras{plumier} \textsuperscript{(31)}. Son livre a beaucoup de
\VUgras{pages}; chaque chapitre contient plusieurs \VUgras{paragraphes} et
chaque \VUgras{ligne}, beaucoup de \VUgras{mots}.

%%K t01-09-4 p
Son voisin Étienne \textsuperscript{(24)} est appliqué. Il note dans son cahier
la leçon pour la prochaine fois. Il a une belle \VUgras{écriture}. Son livre
est fermé. On n'en voit que la \VUgras{couverture} \textsuperscript{(27)}. Son
\VUgras{compas} \textsuperscript{(26)} est près de lui.

%%K t01-09-5 p
Charles \textsuperscript{(23)} tient son livre à la main; il l'ouvre pour
relire sa leçon. Comme chaque élève, il a un \VUgras{encrier}
\textsuperscript{(25)} devant lui. Il est obéissant.

%%K t01-10-1 p
10. --- À la table voisine se trouvent \VUgras{Louis}, \VUgras{Antoine} et
\VUgras{Paul}. Louis récite sa \VUgras{leçon} \textsuperscript{(22)} et répond
aux \VUgras{questions} du professeur. Antoine \textsuperscript{(21)} est très
distrait; le professeur le punit même parfois. Paul \textsuperscript{(20)} est
un bon camarade, aimable et serviable. Il est très attentif.

%%K t01-tit-5 sub
\VUsaut{4.2mm}
\VUpk{10.2pt}{40}{Les mauvais élèves.}
\VUsaut{3.2mm}

%%K t01-11-1 p
11. --- Derrière Henri se trouve \VUgras{Georges} \textsuperscript{(38)}. Ce
n'est pas un mauvais garçon, mais il est \VUgras{bavard}, distrait et
turbulent. Sa \VUgras{légèreté} et sa \VUgras{désobéissance} causent du
\VUgras{désordre} pendant la leçon, et souvent il mérite un sévère
\VUgras{avertissement}. Son \VUgras{cahier de brouillon}
\textsuperscript{(35)} est sale, il le \VUgras{tache d'encre} avec sa
\VUgras{plume} \textsuperscript{(34)}; c'est pourquoi il se sert toujours de
son \VUgras{efface} \textsuperscript{(36)}; sa \VUgras{serviette}
\textsuperscript{(33)} est déchirée, parce qu'il est très négligent. Pourquoi
apporte-t-il son \VUgras{étui à crayons} \textsuperscript{(37)} dans la salle
de classe des \VUgras{langues vivantes}?

%%K t01-11-2 p
Son voisin Edmond \textsuperscript{(39)} ne l'aime pas. Il est vrai qu'il est
un peu paresseux, mais il est silencieux et tranquille. Il ne bavarde pas;
seulement, au lieu d'écouter, il préfère lire les \VUgras{récits} de son
\VUgras{livre de lecture}.

%%K t01-11-3 p
Le troisième élève de ce banc est \VUgras{Louis} \textsuperscript{(40)}. Il est
appliqué. Il écrit sur une \VUgras{feuille de papier} blanche. Devant lui, il a
posé son \VUgras{papier buvard} \textsuperscript{(41)}.

%%K t01-12-1 p
12. --- Les élèves du deuxième banc, à gauche du professeur, sont très jeunes.
Le premier, le petit \VUgras{Émile}, ne sait pas encore très bien écrire; il
apporte son \VUgras{ardoise} \textsuperscript{(42)}, son \VUgras{crayon
d'ardoise} et son \VUgras{éponge} \textsuperscript{(43)}.

%%K t01-12-2 p
À côté de lui se trouvent \VUgras{Auguste} et \VUgras{Bernard}
\textsuperscript{(44)}. Ce dernier travaille bien, mais sa \VUgras{tenue} est
très négligée.

%%K t01-13-1 p
13. --- Derrière eux se trouvent trois \VUgras{paresseux}. \VUgras{Albert}
n'apprend pas ses leçons. \VUgras{Jacques} \textsuperscript{(47)} dort au lieu
d'écouter. Sa \VUgras{paresse} lui vaut des \VUgras{reproches} et même des
\VUgras{punitions}. Enfin, \VUgras{Christian} \textsuperscript{(48)} n'est pas
assez sérieux. Il se \VUgras{conduit} mal et ne fait aucun effort pour se
corriger.

%%K t01-14-1 p
14. --- Ses voisins, au contraire, ont une \VUgras{bonne conduite}.
\VUgras{Ernest} \textsuperscript{(51)} aime l'ordre et la régularité; il se
sert toujours de sa \VUgras{règle} \textsuperscript{(50)} et de son
\VUgras{crayon} \textsuperscript{(49)}. Il est \VUgras{externe}.

%%K t01-14-2 p
\VUgras{François} \textsuperscript{(52)} est \VUgras{pensionnaire}; il porte
l'uniforme, il mange et il couche au lycée. C'est un bon \VUgras{camarade}.

%%K t01-14-3 p
Maurice taille son \VUgras{crayon} avec son \VUgras{canif}
\textsuperscript{(56)}; il est très soigneux.

%%K t01-14-4 p
Il apporte tous ses livres, sa \VUgras{grammaire} \textsuperscript{(55)}, son
\VUgras{carnet de notes} \textsuperscript{(54)}, et même un
\VUgras{sous-main} \textsuperscript{(53)}. Son \VUgras{sac d'école}
\textsuperscript{(57)} est par terre derrière lui.

%%K t01-14-5 p
Les deux tables du fond de la salle de classe sont occupées par de \VUgras{bons
élèves} \textsuperscript{(19)}.

%%K t01-tit-6 sub
\VUsaut{4.6mm}
\VUpk{10.2pt}{40}{Dessin et musique.}
\VUsaut{3.4mm}

%%K t01-15-1 p
15. --- Dans la \VUgras{salle de dessin}, il y a de gracieux \VUgras{modèles}
accrochés au mur et une \VUgras{lampe} \textsuperscript{(62)} avec un
\VUgras{abat-jour}, suspendue au plafond. Le \VUgras{professeur}
\textsuperscript{(60)} est très patient; il corrige les \VUgras{dessins}
\textsuperscript{(61)} des élèves et leur apprend à dessiner un \VUgras{buste}
\textsuperscript{(59)} d'après nature.

%%K t01-16-1 p
16. --- La \VUgras{salle de musique} semble très intéressante. On y apprend à
lire la \VUgras{musique}, à solfier les \VUgras{notes de la gamme}
\textsuperscript{(67)} écrites sur la \VUgras{portée} (do, ré, mi, fa, sol, la,
si\,=\,C. D. E. F. G. A. B.), à reconnaître les \VUgras{clés} et à chanter de
charmantes \VUgras{mélodies}. On pose les partitions sur le \VUgras{lutrin}
\textsuperscript{(65)}. Un des élèves \VUgras{joue du piano}
\textsuperscript{(64)} et le \VUgras{professeur de musique}
\textsuperscript{(66)} \VUgras{bat la mesure} \textsuperscript{(68)}.

%%K t01-17-1 p
17. --- On commence à voir un coin de l'\VUgras{atelier} de \VUgras{travaux
manuels} \textsuperscript{(69)}, où les garçons peuvent apprendre à raboter le
bois et à limer le fer. Cela n'existe pas dans tous les lycées.

%%K t01-tit-7 sub
\VUsaut{5.0mm}
\VUpk{10.2pt}{40}{La salle des sciences.}
\VUsaut{3.6mm}

%%K t01-18-1 p
18. --- Le professeur de mathématiques enseigne aux élèves la
\VUgras{géométrie}, l'\VUgras{arithmétique}, le \VUgras{calcul} et le
\VUgras{système métrique}. On voit sur le tableau les principales figures de
géométrie : le \VUgras{cercle} \textsuperscript{(a)}, le \VUgras{carré}
\textsuperscript{(b)}, le \VUgras{triangle} \textsuperscript{(c)}, la
\VUgras{ligne} \textsuperscript{(d)}.

%%K t01-18-2 p
On voit aussi une \VUgras{opération}; ce n'est pas une \VUgras{addition}, ni
une \VUgras{soustraction}, ni une \VUgras{multiplication} : c'est bien une
\VUgras{division} \textsuperscript{(e)}.

%%K t01-19-1 p
19. --- Sur le tableau du système métrique, on voit : le \VUgras{mètre}
\textsuperscript{(a)}, la \VUgras{balance} \textsuperscript{(b)} et ses
\VUgras{poids}, le \VUgras{litre} \textsuperscript{(e)}, le
\VUgras{décalitre} \textsuperscript{(f)}, le \VUgras{décilitre} et le
\VUgras{mètre cube} \textsuperscript{(d)}.

%%K t01-19-2 p
Les élèves étudient aussi les \VUgras{sciences naturelles}
\textsuperscript{(11)} --- la zoologie \textsuperscript{(a)}, la botanique
\textsuperscript{(b)}, la géologie \textsuperscript{(c)} --- et
l'\VUgras{histoire} \textsuperscript{(12)}.

%%K t01-19-3 p
Dans l'armoire se trouvent les appareils de \VUgras{physique}
\textsuperscript{(15)} et de \VUgras{chimie} \textsuperscript{(16)}.

%%K t01-19-4 p
Sur l'armoire, il y a : une \VUgras{sphère} \textsuperscript{(14)}, un
\VUgras{cône} et un \VUgras{globe terrestre} \textsuperscript{(13)}.

%%K t01-19-5 p
Au-dessus des tableaux est suspendue une \VUgras{carte} qui représente les cinq
parties du monde : l'\VUgras{Europe} \textsuperscript{(g)}, l'\VUgras{Asie}
\textsuperscript{(e)}, l'\VUgras{Afrique} \textsuperscript{(f)},
l'\VUgras{Amérique} \textsuperscript{(ab)} et l'\VUgras{Océanie}
\textsuperscript{(h)}, et les deux grands océans, le \VUgras{Pacifique}
\textsuperscript{(c)} et l'\VUgras{Atlantique} \textsuperscript{(d)}.
\VUsaut{6.0mm}
\VUfilet{22mm}
